\chapter{Cross-platform considerations}
\label{ch:crossplat2}

\begin{aside}
\maya{} runs on Linux, macOS, and BSD. Windows is a different
universe: ConPTY, different signals, different escape parsing.
This chapter is about the things that differ.
\end{aside}

\section{Termios vs Console API}

Unix: \code{tcsetattr} switches to raw mode. Windows:
\code{SetConsoleMode} with \code{ENABLE\_VIRTUAL\_TERMINAL\_
PROCESSING}.

\maya{} wraps both in \code{Terminal::raw\_mode()}.

\section{Signal differences}

\begin{itemize}[leftmargin=*]
  \item SIGWINCH (resize): Unix only. Windows uses a
    different mechanism.
  \item SIGTSTP/SIGCONT (suspend): Unix only.
  \item SIGINT: both have it; Windows via
    \code{SetConsoleCtrlHandler}.
\end{itemize}

\section{Path conventions}

Unix paths: \code{/home/user}. Windows: \code{C:\\Users\\user}.
\code{std::filesystem::path} abstracts but you may still
hit edge cases (case sensitivity, max path length).

\section{Locale defaults}

Unix: \code{en\_US.UTF-8} or similar. Windows: code-page
based; recent versions can use UTF-8.

\section{Terminal differences}

\begin{itemize}[leftmargin=*]
  \item macOS Terminal: missing many features (truecolour,
    DEC 2026, image protocols).
  \item iTerm2 (macOS): rich features.
  \item Windows Terminal: modern; mostly compatible with
    xterm.
  \item Linux console (without X): limited; no truecolour.
\end{itemize}

\section{Fork and exec}

\code{fork} doesn't exist on Windows. Use
\code{CreateProcess}. \maya{}'s spawn API abstracts the
two.

\section{File watching}

Linux: \code{inotify}. macOS: \code{FSEvents} or
\code{kqueue}. BSD: \code{kqueue}. Windows:
\code{ReadDirectoryChangesW}.

\maya{} ships a small abstraction; for advanced use,
\code{libuv}.

\section{Endianness and word size}

\maya{} assumes 64-bit little-endian. Big-endian and 32-bit
are not actively tested. The packed cell type relies on
specific bit layouts.

\section{Recap}

\begin{recap}
\begin{itemize}[leftmargin=*]
  \item Raw mode differs; \maya{} abstracts.
  \item Unix signals don't all map to Windows.
  \item \code{std::filesystem} for paths.
  \item Spawn API differs; abstract.
  \item 64-bit LE assumed.
\end{itemize}
\end{recap}

\section*{Workshop}

\begin{enumerate}[leftmargin=*]
  \item Test your app on a different OS. List the things
    that broke.
  \item Use the spawn abstraction instead of direct fork.
  \item Try a non-x86 build (ARM Mac, Raspberry Pi).
\end{enumerate}
